% =========================================================
% mechlib/elements/bodies.tex
% ATOMARE KOERPER: MASSE, PUNKTMASSE, WAGEN UND BALKEN
% =========================================================
% Die hier definierten TikZ-Pics stellen mechanische Grundkoerper bereit. Jedes
% Pic besitzt eine eigene pgfkeys-Familie unter /mech/... und kann dadurch mit
% benannten Optionen konfiguriert werden.
%
% Grundprinzip der API:
%   \pic (name) at (<Ort>) {mech/<element>={<key>=<wert>,...}};
%
% Benannte Pics exportieren Koordinaten bzw. werden bei node-basierten Koerpern
% auf den internen Shape-Knoten aliasiert. So koennen regulaere TikZ-Anker wie
% (m1.east), (m1.north west) usw. direkt benutzt werden.
%
% Oeffentliche Pics in dieser Datei:
%   mech/mass       rechteckiger Starrkoerper bzw. konzentrierte Masse
%   mech/point-mass kreisfoermige Punktmasse
%   mech/cart       Wagen mit zwei Raedern und optionalem Boden
%   mech/beam       Balken zwischen zwei beliebigen Punkten
% =========================================================
\pgfkeys{
  /mech/mass/.is family,
  /mech/mass,
  width/.initial=1.5cm,
  height/.initial=1cm,
  label/.initial=,
  style/.initial=mech/body,
}
\tikzset{
  pics/mech/mass/.style args={#1}{
    code={
      \pgfkeys{/mech/mass/.cd,#1}
      \node[
        \pgfkeysvalueof{/mech/mass/style},
        minimum width=\pgfkeysvalueof{/mech/mass/width},
        minimum height=\pgfkeysvalueof{/mech/mass/height},
        inner sep=0pt
      ] (-shape) at (0,0) {\pgfkeysvalueof{/mech/mass/label}};

      % Named body pics behave like regular TikZ nodes.
      % Example: \pic (m1) ... then (m1.east), (m1.south east), ...
      \edef\mechPicName{\pgfkeysvalueof{/tikz/name prefix}}
      \ifx\mechPicName\pgfutil@empty\else
        \pgfnodealias{\mechPicName}{\mechPicName-shape}
      \fi

      \coordinate (-center) at (-shape.center);

      % Three standardized connection levels on each vertical side.
      % The shared \mechConnectionOffset is also used by fixed supports.
      \coordinate (-west-upper)  at ($(-shape.west)+(0,\mechConnectionOffset)$);
      \coordinate (-west-middle) at (-shape.west);
      \coordinate (-west-lower)  at ($(-shape.west)+(0,-\mechConnectionOffset)$);

      \coordinate (-east-upper)  at ($(-shape.east)+(0,\mechConnectionOffset)$);
      \coordinate (-east-middle) at (-shape.east);
      \coordinate (-east-lower)  at ($(-shape.east)+(0,-\mechConnectionOffset)$);

      % Generic connection aliases
      \coordinate (-connection-upper)  at (-east-upper);
      \coordinate (-connection-middle) at (-east-middle);
      \coordinate (-connection-lower)  at (-east-lower);

      \coordinate (-north) at (-shape.north);
      \coordinate (-south) at (-shape.south);
    }
  },
  pics/mech/mass/.default={},
}
\pgfkeys{
  /mech/point-mass/.is family,
  /mech/point-mass,
  radius/.initial=2.2mm,
  label/.initial=,
  style/.initial=mech/body,
}
\tikzset{
  pics/mech/point-mass/.style args={#1}{
    code={
      \pgfkeys{/mech/point-mass/.cd,#1}
      \edef\mechPointMassRadius{\pgfkeysvalueof{/mech/point-mass/radius}}
      \edef\mechPointMassLabel{\pgfkeysvalueof{/mech/point-mass/label}}
      \edef\mechPointMassStyle{\pgfkeysvalueof{/mech/point-mass/style}}
      \node[
        \mechPointMassStyle,
        circle,
        minimum size=2\mechPointMassRadius,
        inner sep=0pt
      ] (-shape) at (0,0) {\mechPointMassLabel};

      \edef\mechPicName{\pgfkeysvalueof{/tikz/name prefix}}
      \ifx\mechPicName\pgfutil@empty\else
        \pgfnodealias{\mechPicName}{\mechPicName-shape}
      \fi

      \coordinate (-center) at (-shape.center);
    }
  },
  pics/mech/point-mass/.default={},
}
% -------------------------------------------------------------------------
% Wagen
% -------------------------------------------------------------------------
% Der Boolean "ground" wird als TeX-if gespeichert. Dadurch kann die optionale
% Bodendarstellung ohne Stringvergleich ein- bzw. ausgeschaltet werden.
\newif\ifmechCartGround
\pgfkeys{
  /mech/cart/.is family,
  /mech/cart,
  width/.initial=2cm,
  height/.initial=1cm,
  wheel radius/.initial=2mm,
  label/.initial=,
  ground/.is if=mechCartGround,
  ground/.default=true,
  ground=true,
}
\tikzset{
  pics/mech/cart/.style args={#1}{
    code={
      \pgfkeys{/mech/cart/.cd,#1}
      \edef\mechCartWidth{\pgfkeysvalueof{/mech/cart/width}}
      \edef\mechCartHeight{\pgfkeysvalueof{/mech/cart/height}}
      \edef\mechCartWheelRadius{\pgfkeysvalueof{/mech/cart/wheel radius}}
      \node[
        mech/body,
        minimum width=\mechCartWidth,
        minimum height=\mechCartHeight,
        inner sep=0pt
      ] (-shape) at (0,0) {\pgfkeysvalueof{/mech/cart/label}};

      \edef\mechPicName{\pgfkeysvalueof{/tikz/name prefix}}
      \ifx\mechPicName\pgfutil@empty\else
        \pgfnodealias{\mechPicName}{\mechPicName-shape}
      \fi

      \coordinate (-center) at (-shape.center);
      \coordinate (-west) at (-shape.west);
      \coordinate (-east) at (-shape.east);
      \coordinate (-north) at (-shape.north);
      \coordinate (-south) at (-shape.south);
      \draw[mech/body]
        (-.3\mechCartWidth,-.5\mechCartHeight-\mechCartWheelRadius)
        circle[radius=\mechCartWheelRadius];
      \draw[mech/body]
        (.3\mechCartWidth,-.5\mechCartHeight-\mechCartWheelRadius)
        circle[radius=\mechCartWheelRadius];
      \ifmechCartGround
        \draw
          (-.5\mechCartWidth,-.5\mechCartHeight-2\mechCartWheelRadius)
          --
          (.5\mechCartWidth,-.5\mechCartHeight-2\mechCartWheelRadius);
        \fill[
          pattern={Lines[angle=\mechHatchAngle,distance=1pt]},
          pattern color=black
        ]
          (-.5\mechCartWidth,-.5\mechCartHeight-2\mechCartWheelRadius)
          rectangle
          (.5\mechCartWidth,-.5\mechCartHeight-2\mechCartWheelRadius-3pt);
      \fi
    }
  },
  pics/mech/cart/.default={},
}
\pgfkeys{
  /mech/beam/.is family,
  /mech/beam,
  from/.initial={(0,0)},
  to/.initial={(1,0)},
  height/.initial=3mm,
  style/.initial=mech/body,
}
\tikzset{
  pics/mech/beam/.style args={#1}{
    code={
      \pgfkeys{/mech/beam/.cd,#1}
      \edef\mechBeamFrom{\pgfkeysvalueof{/mech/beam/from}}
      \edef\mechBeamTo{\pgfkeysvalueof{/mech/beam/to}}
      \edef\mechBeamHeight{\pgfkeysvalueof{/mech/beam/height}}
      \edef\mechBeamStyle{\pgfkeysvalueof{/mech/beam/style}}
      \path let
        \p1=\mechBeamFrom,
        \p2=\mechBeamTo,
        \n1={veclen(\x2-\x1,\y2-\y1)},
        \n2={(\y2-\y1)/\n1*\mechBeamHeight/2},
        \n3={-(\x2-\x1)/\n1*\mechBeamHeight/2}
      in
        coordinate (-A) at (\x1+\n2,\y1+\n3)
        coordinate (-B) at (\x2+\n2,\y2+\n3)
        coordinate (-C) at (\x2-\n2,\y2-\n3)
        coordinate (-D) at (\x1-\n2,\y1-\n3)
        coordinate (-start) at (\x1,\y1)
        coordinate (-end) at (\x2,\y2)
        coordinate (-center) at ({(\x1+\x2)/2},{(\y1+\y2)/2});
      \draw[\mechBeamStyle] (-A)--(-B)--(-C)--(-D)--cycle;
    }
  },
  pics/mech/beam/.default={},
}
